\begin{abstract}
Critic-free reinforcement learning has become the standard way to train language-model agents, but it
assigns credit coarsely: GRPO gives every token of a trajectory the same advantage, and GiGPO and
GACA refine this to the level of an environment step. Inside a step, however, a reasoning model
writes a hundred tokens of chain of thought before a ten-token action, and all of them still share one
advantage. We ask where in that response the decision is actually made. A counterfactual flip test on
ALFWorld (truncate the reasoning, resample the rest, and check whether the action changes) shows that
82\% of the decision is settled before the action is written, that token entropy is a poor guide to
where (the final tenth of a response holds 19\% of the decision but 3\% of the entropy), and that a cheap
teacher-forced \emph{commitment probe}, the measurement behind Causal Importance of Reasoning, tracks
it closely: weighting tokens by the probe's rise places $1.65\times$ more credit on decisive segments
than uniform credit, consistently across task types, held-out games and a trained policy, whereas the
model's own attention from action to reasoning does so only weakly ($1.12\times$). We turn the probe into COMMIT (COMMItment-weighted Token credit), a mean-preserving,
clipped redistribution of each step's GACA advantage over its tokens. In a budget-matched study on
ALFWorld with Qwen2.5-1.5B-Instruct and full fine-tuning (about 6\% of GACA's training budget, same
128 held-out games for every method), COMMIT reaches the highest final success of all
arms (20.3\% against 18.8\% for GACA and 14.1\% for GRPO), but none of the method differences is
statistically resolved at this budget, and a variant that applies the probe to positive advantages only
performs significantly worse than both baselines. We also find that
GACA's per-step criticality score is about 85\% within-state noise. Code, logs and weights are released.
\end{abstract}
